\subsection{可测子集上定义的可测函数}

命题 15. --- 设 X 是一个局部紧空间，\(\mu\) 是 X 上的一个测度，A 是 X 的一个 \(\mu\)-可测子集，f 是 A 到一个拓扑空间 F 中的映射。下列条件等价：

\begin{enumerate}
\def\labelenumi{\alph{enumi})}
\item
  集族 \(\mathfrak{H}\) 由这样的紧子集 \(K\)（\(A\) 的）组成：\(f\) 到 \(K\) 上的限制连续；它 \(\mu\)-稠密于 \(A\)。
\item
  存在一个集族 \(\mathfrak{K}\)（由 A 的紧子集组成），它 \(\mu\)-稠密于 A，且 \(f\) 到每个 \(\mathrm{K} \in \mathfrak{K}\) 上的限制是 \(\mu_{\mathrm{K}}\)-可测的。
\item
  存在一个同胚 \(j\)（\(\mathrm{F}\) 到一个拓扑空间 \(\mathrm{G}\) 的一个子空间上的），以及一个 \(\mu\)-可测映射 \(g\)（\(\mathrm{X}\) 到 \(\mathrm{G}\) 中的），使得 \(g|_{\mathrm{A}} = j \circ f\) 。
\item
  f 到 X 上的每个延拓（作为 X 到 F 中的映射，且在 X - A 上取常值）都是 \(\mu\)-可测的。
\end{enumerate}

显然 \(a)\) 蕴涵 \(b)\)，且 \(d)\) 蕴涵 \(c)\) 。\(c)\) 蕴涵 \(a)\) 这一事实由 No.~8 命题 12 的条件 \(c)\) 得出。另一方面，\(b)\) 蕴涵 \(a)\) ：因为，定义 1 表明，对每个 \(K \in \mathfrak{K}\)，子集 \(H \in \mathfrak{H}\)（含于 \(K\) 者）所成之集 \(\mu_{K}\)-稠密（在 \(K\) 中）（No.~8, 命题 12, \(c\) ），而 No.~8 的命题 13 表明 \(\mathfrak{H}\) 是 \(\mu\)-稠密的（在 \(A\) 中）。剩下要证 \(a)\) 蕴涵 \(d)\) 。设 \(g\) 是 \(f\) 到 \(X\) 上的一个延拓，它在 \(X - A\) 上取常值。对每个紧子集 \(L\)（\(X\) 的），\(L \cap A\) 与 \(L \cap (X - A)\) 都是 \(\mu\)-可积的；因此，对每个 \(\varepsilon > 0\)，存在紧子集 \(P \subset L \cap A\) 与紧子集 \(Q \subset L \cap (X - A)\)，使得

\[
| \mu | ((L \cap A) - P) \leqslant \varepsilon / 4 \quad \text { and } \quad | \mu | ((L \cap (X - A)) - Q) \leqslant \varepsilon / 4.
\]

另一方面，存在含于 P 的集 \(H \in \mathfrak{H}\)，使得 \(|\mu|(P - H) \leqslant \varepsilon/2\)；于是 g 到紧集 \(K = H \cup Q\) 上的限制是连续的（g 在 Q 上取常值），且 \(|\mu|(L - K) \leqslant \varepsilon\)，证毕。

定义 8. --- 设 X 是一个局部紧空间，\(\mu\) 是 X 上的一个测度，A 是 X 的一个 \(\mu\)-可测子集。称 A 到一个拓扑空间 F 中的映射 f 是 \(\mu\)-可测的，如果它满足命题 15 的等价条件。

因此，若 A 是局部 \(\mu\)-可忽略的，则 A 到 F 中的每个映射都是 \(\mu\)-可测的。

推论 1. --- 设 X 是一个局部紧空间，\(\mu\) 是 X 上的一个测度，A 是 X 的一个 \(\mu\)-可测子集，\(f\) 是 A 到一个拓扑空间 F 中的一个 \(\mu\)-可测映射。设 \(\mathfrak{K}\) 是 X 的紧子集所成的一个集族，在 X 中 \(\mu\)-稠密。则存在 A 的一个分割，它由一个局部可忽略集 N 与一个局部可数族 \((\mathrm{K}_{\lambda})_{\lambda \in \mathrm{L}}\)（诸集 \(\mathrm{K}_{\lambda} \in \mathfrak{K}\)）组成，使得 \(f|_{\mathrm{K}_{\lambda}}\) 对每个 \(\lambda \in \mathrm{L}\) 连续。

鉴于 No.~9 命题 14，只需证明集族 \(\mathfrak{H} \subset \mathfrak{K}\)——由子集 \(K \in \mathfrak{K}\)（满足 \(K \subset A\) 且 \(f|K\) 连续者）组成——是 \(\mu\)-稠密的（在 \(A\) 中）。而由 No.~1 的命题 1 与条件 \(d)\)（命题 15 的）立即得出：对每个紧子集 \(K_0\)（\(A\) 的） 与每个 \(\varepsilon > 0\)，存在子集 \(K \subset K_0\)（属于 \(\mathfrak{K}\)），使得 \(|\mu|(K_0 - K) \leqslant \varepsilon\) 且 \(f|K\) 连续；因此结论由 No.~8 命题 12 得出。

推论 2. --- 设 K 是 X 的一个紧子空间；为使 K 到一个拓扑空间 F 中的映射 f 是 \(\mu\)-可测的，其充要条件是它是 \(\mu_{K}\)-可测的。

鉴于 No.~7 引理 2，这由 No.~1 命题 1 与命题 15 的条件 \(a\) ) 直接得出。

命题 16. --- 设 A 是 X 的 μ-可测子集所成的一个局部可数集族，令 \(B = \bigcup_{A \in \mathfrak{A}} A\) 。为使 B 到一个拓扑空间 F 中的映射 f 是 μ-可测的，其充要条件是 f 到每个 \(A \in \mathfrak{A}\) 上的限制都是 μ-可测的。

我们已经指出（No.~9），B 是 \(\mu\)-可测的。条件显然必要，我们来证明它充分。于是，设 K 是 B 的一个紧子集。由假设，存在一个集合序列 \((\mathbf{A}_{n})\)（属于 \(\mathfrak{A}\)），使诸 \(K \cap A_{n}\) 构成 K 的一个覆盖。令 \(C_{0} = K \cap A_{0}\)，并令

\[
\mathrm{C} _ {n} = \mathrm{K} \cap \mathrm{A} _ {n} \cap \mathbf {C} \left(\bigcup_ {i <   n} \mathrm{C} _ {i}\right)
\]

（对 n \textgreater{} 0），使诸非空的 \(C_{n}\) 构成 K 的一个由 \(\mu\)-可积集组成的分割。由于 f 到 \(C_{n}\) 上的限制是 \(\mu\)-可测的，存在 \(C_{n}\) 的一个分割，由一个 \(\mu\)-可忽略集 \(N_{n}\) 与一个紧集序列 \((\mathrm{L}_{mn})_{m \geqslant 0}\) 组成，使 \(f|L_{mn}\) 连续。由于 \(N = \bigcup_{n} N_{n}\) 是 \(\mu\)-可忽略的，我们看出命题 15 的条件 a) 满足，命题得证。

命题 15 的性质 d) 使我们可以把 No.~2 至 No.~5 中对定义在整个 X 上的可测函数所观察到的性质，立即推广到定义在 X 的可测子集 A 上的可测函数；这些推广留给读者。我们只明确指出，局部化原理（No.~2, 命题 4）在假设每个函数 \(g_{x}\) 仅在 \(V_{x}\) 中（或局部几乎处处在 \(V_{x}\) 中）有定义且可测时仍然成立。

